\section{AI-Assisted Code Generation}
\label{sec:ai_generative}

Recent advances in generative artificial intelligence have significantly influenced modern software development practices. The introduction of the Transformer architecture established the foundation for large-scale language modeling by replacing recurrent computation with an attention mechanism that captures long-range dependencies more effectively \cite{Vaswani2017Attention}. Building upon this architecture, large language models (LLMs) trained on massive text corpora demonstrated the ability to perform a wide range of language tasks with minimal task-specific supervision \cite{Brown2020GPT3}. When trained on large collections of publicly available source code, these models further acquired the capability to generate syntactically correct and functionally relevant program code from natural language descriptions, enabling a new paradigm of AI-assisted software development \cite{Chen2021Codex}.

Large language models trained on code have since been integrated into practical software development tools, allowing developers to generate code, complete functions, and translate natural language requirements into executable programs. These capabilities extend naturally to infrastructure-as-code (IaC), where cloud infrastructure is defined declaratively through configuration templates or programmatic frameworks. By generating IaC artifacts directly from natural language prompts, AI-assisted code generation can substantially accelerate infrastructure provisioning and reduce the manual effort required to author complex cloud configurations. This capability is particularly valuable in hybrid cloud environments, where infrastructure must be provisioned consistently across heterogeneous platforms with different deployment mechanisms.

Despite these advantages, AI-assisted code generation introduces important security concerns. Because large language models learn statistical patterns from large-scale training data that may contain insecure or outdated coding practices, the generated code does not inherently guarantee security, correctness, or compliance with organizational policies. Empirical studies have shown that a significant proportion of code produced by AI programming assistants contains security vulnerabilities. Pearce et al. evaluated code generated by GitHub Copilot and found that a substantial fraction of the generated programs exhibited security weaknesses corresponding to well-known vulnerability categories \cite{Pearce2022Asleep}. Similarly, controlled user studies have observed that developers assisted by large language models tend to produce less secure code while simultaneously expressing greater confidence in its correctness, suggesting that AI assistance may inadvertently amplify security risks if left unchecked \cite{Perry2023Insecure}.

These concerns are particularly relevant in the context of infrastructure-as-code, where insecure configurations can directly expose cloud environments to attack. AI-generated IaC may introduce misconfigurations such as overly permissive access policies, unencrypted storage resources, or publicly exposed network endpoints, which can propagate throughout the infrastructure once deployed. In hybrid cloud environments, the impact of such misconfigurations is further amplified because infrastructure is distributed across multiple platforms with differing security controls. Consequently, the increased development velocity provided by AI-assisted code generation must be balanced with automated security verification to prevent insecure infrastructure from being deployed.

These observations motivate the integration of AI-assisted code generation with an automated security assessment mechanism within the operational workflow proposed in this research. Rather than treating AI-generated infrastructure code as trustworthy by default, the generated artifacts are systematically evaluated through security scanning and risk assessment before deployment. Furthermore, the security findings identified during evaluation can be provided back to the generation process as feedback, enabling iterative refinement of the generated code. By coupling AI-assisted code generation with continuous security verification, the proposed approach aims to leverage the productivity benefits of generative AI while mitigating the security risks associated with automatically generated infrastructure-as-code.
